{\gostTitleFont
    \redline
    СОДЕРЖАНИЕ
} 

\titlespace
 
{\gostFont
    \par
   \par {\bfseries ВВЕДЕНИЕ \dotfill 5}
   \par {\bfseries 1. АНАЛИЗ ПРЕДМЕТНОЙ ОБЛАСТИ И ПОСТАНОВКА ЗАДАЧИ \dotfill 6}
   \par 1.1 Общие сведения о сервисах заказа такси \dotfill 6
   \par 1.2 Особенности разработки высоконагруженных серверных систем \dotfill 8
   \par 1.3 Сравнительный анализ монолитной и микросервисной архитектур \dotfill 10
   \par 1.4 Требования к серверной части приложения и постановка задачи \dotfill 12
   \par 1.5 Анализ существующих решений на рынке такси-сервисов \dotfill 14
   \par {\bfseries 2. ВЫБОР ИСПОЛЬЗУЕМЫХ ТЕХНОЛОГИЙ \dotfill 16}
   \par 2.1 Обзор языков программирования и платформ для серверной разработки \dotfill 16
   \par 2.2 Выбор и обоснование стека технологий (Java, Spring Boot) \dotfill 17
   \par 2.3 Анализ систем управления базами данных и средств кэширования \dotfill 18
   \par {\bfseries 3. ПРОЕКТИРОВАНИЕ АРХИТЕКТУРЫ СЕРВЕРНОЙ ЧАСТИ \dotfill 20}
   \par 3.1 Разработка схемы взаимодействия микросервисов \dotfill 20
   \par 3.2 Проектирование структуры базы данных и геоинформационных индексов \dotfill 22
   \par 3.3 Разработка API и протоколов обмена данными \dotfill 25
   \par {\bfseries 4. РЕАЛИЗАЦИЯ ОСНОВНЫХ МОДУЛЕЙ СИСТЕМЫ \dotfill 28}
   \par 4.1 Реализация сервиса управления заказами \dotfill 28
   \par 4.2 Разработка модуля обработки геолокации и поиска ближайших водителей \dotfill 31
   \par 4.3 Система авторизации и аутентификации на основе JWT \dotfill 33
   \par 4.4 Интеграция брокера сообщений для обработки событий в реальном времени \dotfill 35
   \par {\bfseries 5. ТЕСТИРОВАНИЕ И РАЗВЕРТЫВАНИЕ СИСТЕМЫ \dotfill 39}
   \par 5.1 Модульное и интеграционное тестирование компонентов \dotfill 39
   \par 5.2 Контейнеризация приложения с использованием Docker \dotfill 40
   \par 5.3 Настройка CI/CD процессов и мониторинг серверной части \dotfill 42
   \par {\bfseries 6. ТЕХНИКО-ЭКОНОМИЧЕСКОЕ ОБОСНОВАНИЕ \dotfill 44}
   \par 6.1 Исходные данные для расчета экономического эффекта \dotfill 44
   \par 6.2 Расчёт объёма функций программного продукта \dotfill 44
   \par 6.3 Расчёт полной себестоимости разработки \dotfill 45
   \par 6.4 Оценка экономической эффективности внедрения системы \dotfill 49 
   \par {\bfseries ЗАКЛЮЧЕНИЕ \dotfill 52}
    \par {\bfseries СПИСОК СОКРАЩЕНИЙ \dotfill 53} 
    \par {\bfseries СПИСОК ЛИТЕРАТУРЫ \dotfill 54} 
    \par {\bfseries ПРИЛОЖЕНИЕ А: ТЕКСТ ПРОГРАММЫ \dotfill} 
    \par 
}
